\documentclass[../main.tex]{subfiles}
\begin{document}

\chapter{Memory}
\lessoninfo{
  video={Lesson 15, videos 1--14},
  textbook={H\&P \cite{hennessy2017} §2.2},
  keywords={SRAM, DRAM, memory cell, word line, bit line, sense amplifier, destructive read, trench cell, refresh, row decoder, RAS and CAS, fast page mode, front-side bus, northbridge},
}

Lessons 12--14 treated main memory as a large, slow device hidden behind the
caches, and Lesson 1 showed that the gap between processor and memory speed,
the memory wall, keeps widening.  This lesson looks inside the memory itself.
It describes the two technologies from which memories are built, static RAM
for caches and dynamic RAM for main memory, down to the circuit that stores a
single bit; it then shows how DRAM cells are organized into a chip and
accessed, and how the chip is connected to the processor.  Together these
explain why main memory is slow and why it is not simply built with the
faster technology of the caches.

\section{Memory technologies}
\label{sec:l15-technologies}

Caches and main memory are both random-access memories.\source{Lesson 15,
videos 1--3; H\&P §2.2.}

\begin{definition}[Random access]
  A memory provides \term{random access}[ランダムアクセス] if any location can
  be read or written directly through its address, in about the same time for
  every location.
\end{definition}

The opposite is sequential access (Lesson 16), in which a medium such as a
magnetic tape has to be moved past all the data that lies between the current
position and the requested one.  Random-access memory (RAM) is built in two
technologies, named after what happens to a stored bit while the memory is
powered.

\begin{definition}[SRAM and DRAM]
  In \term{static random-access memory}[SRAM]%
  \index{SRAM|see{static random-access memory}} (SRAM) a stored bit keeps its
  value for as long as the memory is powered.  In
  \term{dynamic random-access memory}[DRAM]%
  \index{DRAM|see{dynamic random-access memory}} (DRAM) a stored bit loses
  its value after a short time, even with power on, unless it is
  periodically rewritten.
\end{definition}

The names say nothing about what happens without power: both technologies are
volatile and lose their entire contents when power is removed.

The two technologies differ in the circuit that stores a bit, and this
difference determines their speed and cost.  SRAM uses several transistors per
bit, six in the usual design (\cref{sec:l15-sram}); DRAM uses a single
transistor (\cref{sec:l15-dram}).  SRAM therefore takes more chip area per bit
and is more expensive per bit, but it is faster, because its cell actively
drives the bit lines, whereas a DRAM cell only holds a small charge that has to
be sensed, written back and refreshed.  DRAM packs more bits into the same area
and is cheaper per bit, but slower.  Caches, which must be fast
and can be small, are built from SRAM; main memory, which must be large and
cheap, is built from DRAM.\margin{The division is not absolute: IBM's POWER7
(2010) built its 32~MB L3 cache from DRAM cells---trench eDRAM---on the
processor die.}

\section{The SRAM cell}
\label{sec:l15-sram}

A memory stores its bits in a two-dimensional array of identical
circuits.\source{Lesson 15, video 4; H\&P §2.2.}  The circuit that stores one
bit is called a \term{memory cell}[メモリセル].  The cells are arranged in
rows and columns and are reached through two sets of wires that cross the
array.

\begin{definition}[Bit line and word line]
  A \term{bit line}[ビット線] runs along a column of cells and carries the
  value that is read from or written to a cell of the column.  A
  \term{word line}[ワード線] runs along a row of cells; when it is activated,
  it connects every cell of the row to the bit line of its column.
\end{definition}

In an SRAM cell (\cref{fig:l15-cells}a) the bit is held by two inverters
connected in a loop, the output of each driving the input of the other.  If
the first inverter outputs~1, the second outputs~0, which in turn keeps the
first at~1.  The loop is stable in either of its two states and holds it for
as long as the inverters are powered, which is what makes SRAM static.  An
access transistor, whose gate is connected to the word line, connects one node
of the loop to the bit line.

\begin{figure}[htbp]
  \centering
  % One memory bit.  (a) 6T SRAM cell: two cross-coupled inverters hold the
% bit; two access transistors, gated by the word line, connect the two
% nodes of the loop to the bit lines B and not-B.  (b) 1T DRAM cell: one
% access transistor connects the bit line to a storage capacitor.
\begin{tikzpicture}[x=1cm,y=1cm, font=\footnotesize\sffamily, line cap=round,
  dot/.style={circle, fill, inner sep=0pt, minimum size=3.2pt},
  lab/.style={font=\scriptsize\sffamily, text=ln-muted}]
  % \lnnmos{x}{y}: horizontal access transistor, channel centred at x, leads
  % at height y, gate wire going up to the word line at y=2.4
  \def\lnnmos#1#2{%
    \draw (#1-0.3,#2) -- (#1-0.3,#2+0.15) -- (#1+0.3,#2+0.15) -- (#1+0.3,#2);
    \draw[very thick] (#1-0.3,#2+0.3) -- (#1+0.3,#2+0.3);
    \draw (#1,#2+0.3) -- (#1,2.4) node[dot] {};
  }
  % ================================================= (a) SRAM cell
  \draw[ln-accent, thick] (0,2.4) -- (6.0,2.4);
  \node[lab, anchor=south] at (3.0,2.42) {\iflangja{ワード線}{word line}};
  \draw[ln-accent, thick] (0.4,2.8) -- (0.4,-1.0) node[below, text=black] {$B$};
  \draw[ln-accent, thick] (5.6,2.8) -- (5.6,-1.0) node[below, text=black] {$\overline{B}$};
  % access transistors
  \lnnmos{1.25}{1.4}
  \draw (0.4,1.4) node[dot] {} -- (0.95,1.4);
  \draw (1.55,1.4) -- (2.2,1.4);
  \lnnmos{4.75}{1.4}
  \draw (3.8,1.4) -- (4.45,1.4);
  \draw (5.05,1.4) -- (5.6,1.4) node[dot] {};
  % storage nodes
  \draw (2.2,1.4) -- (2.2,-0.2) -- (2.8,-0.2);
  \draw (3.8,1.4) -- (3.8,-0.2) -- (3.25,-0.2);
  % upper inverter (left to right)
  \draw[fill=ln-box] (2.75,0.45) -- (2.75,0.95) -- (3.2,0.7) -- cycle;
  \draw[fill=white] (3.26,0.7) circle (0.06);
  \draw (2.2,0.7) node[dot] {} -- (2.75,0.7);
  \draw (3.32,0.7) -- (3.8,0.7) node[dot] {};
  % lower inverter (right to left)
  \draw[fill=ln-box] (3.25,-0.45) -- (3.25,0.05) -- (2.8,-0.2) -- cycle;
  \draw[fill=white] (2.74,-0.2) circle (0.06);
  \draw (2.68,-0.2) -- (2.2,-0.2);
  \node[lab] at (3.0,1.05) {\iflangja{インバータ対}{inverter loop}};
  \node[anchor=north] at (3.0,-1.45) {\iflangja{(a) SRAM セル（6 トランジスタ）}{(a) SRAM cell (6 transistors)}};
  % ================================================= (b) DRAM cell
  \draw[ln-accent, thick] (7.2,2.4) -- (10.6,2.4);
  \node[lab, anchor=south] at (9.7,2.42) {\iflangja{ワード線}{word line}};
  \draw[ln-accent, thick] (7.8,2.8) -- (7.8,-1.0) node[below, text=black, lab] {\iflangja{ビット線}{bit line}};
  \lnnmos{8.65}{1.4}
  \draw (7.8,1.4) node[dot] {} -- (8.35,1.4);
  \draw (8.95,1.4) -- (9.6,1.4) -- (9.6,0.55);
  % capacitor and ground
  \draw[very thick] (9.3,0.55) -- (9.9,0.55);
  \draw[very thick] (9.3,0.38) -- (9.9,0.38);
  \draw (9.6,0.38) -- (9.6,-0.2);
  \draw (9.38,-0.2) -- (9.82,-0.2);
  \draw (9.46,-0.3) -- (9.74,-0.3);
  \draw (9.54,-0.4) -- (9.66,-0.4);
  \node[lab, anchor=west] at (9.95,0.47) {\iflangja{キャパシタ}{capacitor}};
  \node[anchor=north] at (8.9,-1.45) {\iflangja{(b) DRAM セル（1 トランジスタ）}{(b) DRAM cell (1 transistor)}};
\end{tikzpicture}

  \caption{One memory bit.  (a)~An SRAM cell holds the bit in a loop of two
    inverters, whose nodes are connected to the complementary bit lines $B$
    and $\overline{B}$ by two access transistors.  (b)~A DRAM cell holds the
    bit as charge on a capacitor behind a single access transistor.}
  \label{fig:l15-cells}
\end{figure}

\begin{description}
  \item[Write] The bit line is driven strongly to the new value and the word
    line is activated.  The bit line overpowers the small inverter that drives
    the node, the loop settles in the state that matches the bit line, and it
    keeps this state after the word line has been deactivated.
  \item[Read] The word line is activated and the loop drives the bit line.
    The transistors of the inverters are small, because the cell has to be
    small, whereas the bit line is a long wire attached to every cell of the
    column.  A single cell can change the voltage of the bit line only slowly.
\end{description}

To speed up the read, the other node of the loop, which holds the complement
of the bit, is connected through a second access transistor to a second bit
line, $\overline{B}$.  When the word line is activated, the cell pulls the two
bit lines in opposite directions, and the value is found by comparing them.

\begin{definition}[Sense amplifier]
  A \term{sense amplifier}[センスアンプ] detects a small voltage change on a
  bit line, or a small difference between a pair of bit lines, and amplifies
  it to a full logic~0 or~1.
\end{definition}

With a complementary pair, the sense amplifier only compares the two bit
lines.  The cell moves them in opposite directions, so the sign of their
difference is clear after a very small swing and no reference voltage is
needed; a single bit line would have to move much further before its value
could be decided reliably.  The pair helps writes as well,
since the value and its complement act on both sides of the loop at once.  The
resulting cell has six transistors, two in each inverter and two access
transistors, and it needs a word line, two bit lines and the supply
connections of its inverters.\margin{Every cache array of Lessons 12--14 is
built from these cells: the data of a 32~KB L1 cache alone is \num{262144}
bits, some 1.6~million transistors.}

\section{The DRAM cell}
\label{sec:l15-dram}

A DRAM cell stores its bit as electric charge on a capacitor
(\cref{fig:l15-cells}b).\source{Lesson 15, videos 5--6; H\&P §2.2;
\cite{dennard1968}.}  A single access transistor, controlled by the word line,
connects the capacitor to the bit line; a charged capacitor represents~1 and a
discharged one~0.  To write, the bit line is set to the value and the word
line is activated, so that the capacitor charges or discharges through the
transistor.  When the word line is deactivated, the transistor turns off and
isolates the charge.

A transistor that is off is not a perfect switch, however.\caution{The
\enquote{static} of SRAM is about retention, not power: SRAM cells leak as
well, and the caches are among the largest contributors to a processor's
static power.}  A small current
still flows through it---the same imperfection that causes static power
(Lesson 1)---and the charge of a stored~1 gradually leaks away until the cell
can no longer be distinguished from a~0.  This is what makes DRAM dynamic: a
bit survives only if it is rewritten before too much of its charge has been
lost.

\begin{definition}[Refresh]
  Reading the bits of DRAM cells and writing them back with full charge
  before the charge has leaked away is called \term{refresh}[リフレッシュ].
  Every cell must be refreshed periodically, whether or not its data is being
  used.
\end{definition}

Reading a DRAM cell is delicate as well.  The bit line is first set to a
voltage halfway between~0 and~1, and then the word line is activated.  The
capacitor shares its charge with the bit line, whose voltage rises slightly if
the cell held~1 and falls slightly if it held~0.  The capacitor of a cell is
tiny compared with the capacitance of the long bit line, so the change is very
small, and a sense amplifier is needed to detect it.  Once the charge has been
shared, the capacitor is left near the intermediate voltage: the stored value
is gone.

\begin{definition}[Destructive read]
  A \term{destructive read}[破壊読み出し] is a read that destroys the stored
  value.  Every read of a DRAM cell is therefore followed by writing the value
  that was read back into the cell.
\end{definition}

An SRAM read is not destructive, because the inverter loop keeps driving its
nodes and restores their values by itself.

The capacitor also raises a problem of area.  Built flat on the surface of the
chip, a capacitor large enough to be sensed reliably would occupy far more
area than the transistor, and much of the density advantage of a
one-transistor cell would be lost.

\begin{definition}[Trench cell]
  A \term{trench cell}[トレンチセル] is a DRAM cell whose capacitor is formed
  in a deep, narrow trench etched vertically into the silicon next to the
  access transistor (\cref{fig:l15-trench}).
\end{definition}

\begin{marginfigure}
  % Cross-section of a DRAM trench cell (schematic): the access transistor
% sits at the surface; the storage capacitor is a deep, narrow trench
% etched into the silicon next to it.
\resizebox{\linewidth}{!}{%
\begin{tikzpicture}[x=1cm,y=1cm, font=\footnotesize\sffamily,
  lab/.style={font=\footnotesize\sffamily, text=ln-muted}]
  % substrate
  \fill[ln-box] (0,0) rectangle (4.2,-3.2);
  \draw[ln-rule] (0,0) -- (4.2,0);
  % source/drain regions
  \draw[fill=ln-accent!30] (0.5,0) rectangle (1.3,-0.35);
  \draw[fill=ln-accent!30] (2.1,0) rectangle (3.5,-0.35);
  % gate (word line)
  \draw[fill=ln-accent] (1.3,0.08) rectangle (2.1,0.36);
  \draw (1.7,0.36) -- (1.7,0.9);
  \node[lab, anchor=south] at (1.7,0.9) {\iflangja{ワード線}{word line}};
  % bit line contact
  \draw (0.9,0) -- (0.9,0.6) -- (0.2,0.6);
  \node[lab, anchor=south west] at (0.0,0.62) {\iflangja{ビット線}{bit line}};
  % trench capacitor: storage electrode joined to the drain region, lined
  % with a thin dielectric
  \fill[ln-accent!30] (2.68,-0.3) rectangle (3.32,-2.82);
  \draw[ln-muted, line width=1.2pt] (2.64,-0.35) -- (2.64,-2.86) -- (3.36,-2.86) -- (3.36,-0.35);
  \node[rotate=-90, font=\footnotesize\sffamily] at (3.0,-1.6) {\iflangja{容量}{capacitor}};
  \node[lab, anchor=west] at (3.5,-2.55) {Si};
\end{tikzpicture}}

  \caption{Cross-section of a trench cell (schematic).}
  \label{fig:l15-trench}
\end{marginfigure}

The trench gives the capacitor a large plate area while using little of the
chip surface.  The other common design, used in most current DRAMs, stacks the
capacitor above the access transistor; either way the cell occupies little
more area than the transistor.  It also needs fewer wires than an SRAM cell: one bit line instead
of two, and no supply connections for inverters.  \Cref{tab:l15-sram-dram}
summarizes the comparison of the two cells.

\begin{table}[htbp]
  \centering\small
  \caption{SRAM and DRAM compared.}
  \label{tab:l15-sram-dram}
  \begin{tabular}{@{}lll@{}}
    \toprule
                              & SRAM           & DRAM \\
    \midrule
    Storage element           & inverter loop  & capacitor \\
    Transistors per bit       & 6              & 1 \\
    Bit lines per cell        & 2              & 1 \\
    Keeps its value when powered & yes         & only if refreshed \\
    Read                      & non-destructive & destructive, then write-back \\
    Speed                     & fast           & slow \\
    Area and cost per bit     & high           & low \\
    Used for                  & caches         & main memory \\
    \bottomrule
  \end{tabular}
\end{table}

\begin{keypoint}
  DRAM stores a bit as a tiny charge behind one transistor, which makes it
  dense and cheap but slow: every read destroys the charge and must be
  followed by a write-back, and leakage forces a periodic refresh.  SRAM
  avoids both at the price of six transistors per bit.
\end{keypoint}

\section{Memory chip organization}
\label{sec:l15-chip}

The cells of a memory chip form an array of $2^r$ rows and $2^c$ columns
(\cref{fig:l15-chip}), and an address of $r+c$ bits is split into a row
address of $r$ bits and a column address of $c$ bits, for instance the upper
and the lower bits of the address.\source{Lesson 15, video 7; H\&P §2.2.}  A
read of one bit of a DRAM chip proceeds in four steps.
\begin{enumerate}
  \item The \term{row decoder}[行デコーダ] turns the row address into the
    activation of exactly one of the $2^r$ word lines.
  \item Every cell of the selected row is connected to its bit line.  One sense
    amplifier per bit line detects the values, and the whole row is stored in
    a row latch.
  \item The \term{column multiplexer}[列マルチプレクサ] uses the column address
    to select one bit of the latch and delivers it to the data output.
  \item The read was destructive for every cell of the row, so the contents of
    the latch are written back into the row over the bit lines.
\end{enumerate}

\begin{figure}[htbp]
  \centering
  % Organization of a DRAM chip: the row decoder activates one word line
% (row address or refresh counter), the sense amplifiers read the whole row
% into the row latch, the column multiplexer selects one bit, and the latch
% is written back to the row over the bit lines.
\begin{tikzpicture}[x=1cm,y=1cm, font=\footnotesize\sffamily, >={Stealth[length=1.8mm]},
  lab/.style={font=\scriptsize\sffamily, text=ln-muted},
  blk/.style={draw, fill=ln-box}]
  % ---- cell array: one cell at every crossing of a word line and a bit line
  \fill[ln-accent!15] (3.3,2.55) rectangle (9.3,3.05);
  \foreach \i in {0,...,5} {
    \draw[ln-accent] (3.3,{4.0-0.6*\i}) -- (9.3,{4.0-0.6*\i});
  }
  \foreach \j in {0,...,7} {
    \draw[ln-muted] ({3.9+0.7*\j},4.4) -- ({3.9+0.7*\j},0.3);
    \foreach \i in {0,...,5} {
      \draw[fill=white] ({3.9+0.7*\j-0.11},{4.0-0.6*\i-0.11}) rectangle ({3.9+0.7*\j+0.11},{4.0-0.6*\i+0.11});
    }
  }
  \node[lab, anchor=south] at (6.35,4.4) {\iflangja{メモリセルの配列}{array of memory cells}};
  \node[lab, anchor=west] at (9.35,4.0) {\iflangja{ワード線}{word line}};
  \node[lab, anchor=south] at (8.8,4.4) {\iflangja{ビット線}{bit line}};
  % ---- row decoder with address / refresh counter selection
  \draw[blk] (2.6,0.6) rectangle (3.3,4.4);
  \node[rotate=90] at (2.95,2.5) {\iflangja{行デコーダ}{row decoder}};
  \draw[fill=white] (1.7,1.75) -- (2.2,2.05) -- (2.2,2.95) -- (1.7,3.25) -- cycle;
  \draw[->] (2.2,2.5) -- (2.6,2.5);
  \draw[->] (0,2.9) node[above right, inner sep=1pt] {\iflangja{行アドレス}{row address}} -- (1.7,2.9);
  \node[blk, align=center, font=\scriptsize\sffamily, inner sep=2pt] (cnt) at (0.65,1.6)
    {\iflangja{リフレッシュ}{refresh}\\\iflangja{カウンタ}{counter}};
  \draw[->] (cnt.east) -- (1.45,1.6) -- (1.45,2.1) -- (1.7,2.1);
  % ---- sense amplifiers and row latch
  \draw[blk] (3.6,-0.2) rectangle (9.2,0.3);
  \node at (6.4,0.05) {\iflangja{センスアンプ}{sense amplifiers}};
  \draw[blk] (3.6,-1.3) rectangle (9.2,-0.8);
  \node at (6.4,-1.05) {\iflangja{行ラッチ}{row latch}};
  % read into the latch and write back over the same bit lines
  \foreach \j in {0,...,7} {
    \draw[<->, >={Stealth[length=1.5mm]}] ({3.9+0.7*\j},-0.2) -- ({3.9+0.7*\j},-0.8);
  }
  \node[lab, anchor=west, align=left] at (9.25,-0.5)
    {\iflangja{読み出し／}{read /}\\\iflangja{書き戻し}{write back}};
  % ---- column multiplexer
  \foreach \j in {0,...,7} {
    \draw ({3.9+0.7*\j},-1.3) -- ({3.9+0.7*\j},-1.6);
  }
  \draw[fill=white] (3.7,-1.6) -- (9.1,-1.6) -- (7.15,-2.3) -- (5.65,-2.3) -- cycle;
  \node at (6.4,-1.92) {\iflangja{列マルチプレクサ}{column multiplexer}};
  \draw[->] (0.8,-1.95) node[above right, inner sep=1pt] {\iflangja{列アドレス}{column address}} -- (4.68,-1.95);
  \draw[->] (6.4,-2.3) -- (6.4,-2.8) node[below] {\iflangja{データ}{data}};
\end{tikzpicture}

  \caption{Organization of a DRAM chip.  The row decoder activates the word
    line selected by the row address or by the refresh counter
    (\cref{sec:l15-refresh}); the sense amplifiers read the whole row into the
    row latch, the column multiplexer selects one bit, and the latch is written
    back into the row.}
  \label{fig:l15-chip}
\end{figure}

A write uses the same path.  The row is read into the latch as above, the bit
selected by the column address is replaced in the latch by the value to be
written, and the write-back then stores the modified row.  In a DRAM chip a
write is thus a read whose latch is changed before it is written
back.\margin{Real devices are 4, 8 or 16 bits wide rather than one: an
8-bit-wide DDR4 chip holds 1~KB in a row, and the eight chips of a rank open
their rows together (JEDEC DDR4).}

\begin{note}
  \label{note:l15-model}
  A row latch that is written back when the row is closed is the simplified
  model used throughout this lesson.  In real DRAM chips the sense amplifiers
  themselves latch the row: while the word line is still active they drive
  the bit lines to full levels, which restores the cells of the row, and a bit
  written into the latch reaches its cell at once.  Closing the row turns the
  word line off and prepares the bit lines for the next read.
\end{note}

The two parts of the address take very different amounts of time.  Selecting a
row involves decoding, driving a long word line and waiting for the sense
amplifiers to detect tiny charges on long bit lines.  Selecting a column only
switches a multiplexer between bits that are already in the latch.
\Cref{sec:l15-ras-cas,sec:l15-fpm} exploit this difference.\margin{Datasheets
call the latched row the row buffer, and an access that finds its row already
open a \enquote{page hit} (\cref{sec:l15-fpm}).}

\begin{note}
  An SRAM chip is organized in the same way, with a row decoder, sense
  amplifiers and a column multiplexer.  Its reads are not destructive,
  however, so no write-back follows a read, and its cells need no refresh.
\end{note}

\section{Refresh}
\label{sec:l15-refresh}

Every row of a DRAM chip has to be refreshed at least once in every
interval~$T$, typically a few tens of milliseconds.\source{Lesson 15, videos
8--9; H\&P §2.2.}  A refresh uses the normal read path: it is a read
(\cref{sec:l15-chip}) whose latch is written back without using the
output.\margin{JEDEC DDR4 fits \num{8192} refresh commands into a
\SI{64}{\milli\second} window, one every \SI{7.8}{\micro\second}, and halves
the interval above \SI{85}{\celsius}.}

A read by the program refreshes the row it touches, but such reads cannot be
relied upon, because some rows are not read for long periods or not at all.
All rows are therefore refreshed systematically.  A refresh counter holds a
row number.  Periodically the row decoder is given the value of the counter
instead of the row address (\cref{fig:l15-chip}); that row is read and written
back, and the counter advances to the next row, wrapping around after the last
one.  With $R$ rows, each of which must be refreshed once per interval~$T$,
the chip must refresh one row, and advance the counter, at least once in every
interval of length
\begin{equation}
  t_{\text{row}} = \frac{T}{R} .
  \label{eq:l15-refresh}
\end{equation}

While a row is being refreshed, the chip cannot serve a read or a write, and an
access that arrives in the meantime has to wait.  Refresh is thus one more
reason why DRAM is slow, and its cost grows with the number of rows.\tip{The
refresh overhead is $t/(T/R) = Rt/T$, with $t$ the time one row refresh
occupies the chip; how the refreshes are spread over $T$ does not matter.}

\begin{example}
  \label{ex:l15-refresh}
  A DRAM chip has \num{8192} rows whose cells must be refreshed every
  \SI{64}{\milli\second}.  By \cref{eq:l15-refresh} the chip must refresh one
  row in every $\SI{64}{\milli\second}/\num{8192} \approx
  \SI{7.8}{\micro\second}$, i.e.\ \num{128000} rows per second.  If
  refreshing one row takes \SI{50}{\nano\second}, the chip spends
  $\num{128000} \times \SI{50}{\nano\second} = \SI{6.4}{\milli\second}$ of
  every second, or \SI{0.64}{\percent} of the time, on refresh.  A chip with
  four times as many rows and the same~$T$ spends \SI{2.56}{\percent}.
\end{example}

\section{Address multiplexing}
\label{sec:l15-ras-cas}

A chip with $2^{r+c}$ cells needs $r+c$ address bits, and every bit that is
sent in parallel costs a pin on the package and a wire to the
chip.\source{Lesson 15, video 10; H\&P §2.2.}  The row address and the column
address are not needed at the same time, however.  The row address is needed
first, to start the slow selection of the row; the column address is needed
only once the row is in the latch.  DRAM chips therefore receive both parts
over the same pins, one after the other, and two control signals tell the chip
which part is present.

\begin{definition}[RAS and CAS]
  When the \term{row address strobe}[行アドレスストローブ] (RAS) is asserted, the chip takes the
  value on the address pins as the row address, activates the word line and
  reads the row into the latch.  When the
  \term{column address strobe}[列アドレスストローブ]%
  \index{CAS|see{column address strobe}} (CAS) is asserted, the chip takes the
  value on the address pins as the column address and selects that bit of the
  latch.
\end{definition}

On a read the selected bit appears at the data output.  If a write-enable
signal is asserted together with CAS, the bit in the latch is instead replaced
by the value at the data input.  Deasserting RAS ends the access: the latch is
written back and the row is closed (in the model of \cref{note:l15-model}).
\Cref{fig:l15-timing}a shows two
accesses of this kind.\margin{In real chips RAS and CAS are active low;
\cref{fig:l15-timing} draws asserted signals high.}

\begin{figure}[htbp]
  \centering
  % RAS/CAS timing on a DRAM chip with multiplexed address pins (asserted
% signals drawn high).  (a) Two separate accesses to the same row: each
% opens the row with RAS, selects a column with CAS and closes the row.
% (b) Fast page mode: the row is opened once and three columns are
% accessed before RAS is deasserted and the row is written back.
\begin{tikzpicture}[x=1cm,y=1cm, font=\footnotesize\sffamily,
  lab/.style={anchor=east, font=\footnotesize\sffamily},
  ann/.style={font=\scriptsize\sffamily, text=ln-muted}]
  % \lnbus{x0}{x1}{y}{text}: valid value on the address pins
  \def\lnbus#1#2#3#4{%
    \draw[fill=ln-accent!15] (#1,#3) -- (#1+0.1,#3+0.2) -- (#2-0.1,#3+0.2) -- (#2,#3)
      -- (#2-0.1,#3-0.2) -- (#1+0.1,#3-0.2) -- cycle;
    \node[font=\scriptsize\sffamily] at ({(#1+#2)/2},#3) {#4};
  }
  % \lnidle{x0}{x1}{y}: no value on the address pins
  \def\lnidle#1#2#3{\draw[ln-muted] (#1,#3) -- (#2,#3);}
  % \lnsig{y}{list of x0/x1 pulses}: control signal, low outside the pulses
  \def\lnsig#1#2{%
    \draw (1.4,#1-0.2) \foreach \a/\b in {#2} { -- (\a,#1-0.2) -- (\a+0.08,#1+0.2)
      -- (\b-0.08,#1+0.2) -- (\b,#1-0.2) } -- (11.4,#1-0.2);
  }
  % ================================================= (a)
  \node[anchor=west] at (0,0.75) {\iflangja{(a) 同じ行への 2 回のアクセス}{(a) two accesses to the same row}};
  \node[lab] at (1.3,0) {\iflangja{アドレス}{address}};
  \lnidle{1.4}{2.0}{0}
  \lnbus{2.0}{3.4}{0}{\iflangja{行}{row}}
  \lnbus{3.4}{4.8}{0}{\iflangja{列}{column}}
  \lnidle{4.8}{6.6}{0}
  \lnbus{6.6}{8.0}{0}{\iflangja{行}{row}}
  \lnbus{8.0}{9.4}{0}{\iflangja{列}{column}}
  \lnidle{9.4}{11.4}{0}
  \node[lab] at (1.3,-0.6) {RAS};
  \lnsig{-0.6}{2.3/5.1, 6.9/9.7}
  \node[lab] at (1.3,-1.2) {CAS};
  \lnsig{-1.2}{3.7/4.7, 8.3/9.3}
  \node[ann, anchor=west] at (5.15,-0.6) {\iflangja{書き戻し}{write back}};
  \node[ann, anchor=west] at (9.75,-0.6) {\iflangja{書き戻し}{write back}};
  % ================================================= (b)
  \node[anchor=west] at (0,-2.05) {\iflangja{(b) ファストページモード}{(b) fast page mode}};
  \node[lab] at (1.3,-2.8) {\iflangja{アドレス}{address}};
  \lnidle{1.4}{2.0}{-2.8}
  \lnbus{2.0}{3.4}{-2.8}{\iflangja{行}{row}}
  \lnbus{3.4}{4.8}{-2.8}{\iflangja{列}{column}}
  \lnbus{4.8}{6.2}{-2.8}{\iflangja{列}{column}}
  \lnbus{6.2}{7.6}{-2.8}{\iflangja{列}{column}}
  \lnidle{7.6}{11.4}{-2.8}
  \node[lab] at (1.3,-3.4) {RAS};
  \lnsig{-3.4}{2.3/8.2}
  \node[lab] at (1.3,-4.0) {CAS};
  \lnsig{-4.0}{3.7/4.7, 5.1/6.1, 6.5/7.5}
  \node[ann, anchor=west] at (8.2,-3.4) {\iflangja{書き戻し}{write back}};
\end{tikzpicture}

  \caption{Signals of a DRAM chip with multiplexed address pins.  (a)~Each
    access opens the row with RAS, selects a column with CAS and closes the
    row.  (b)~In fast page mode (\cref{sec:l15-fpm}) the row is opened once,
    three columns are accessed, and the row is written back when RAS is
    deasserted.}
  \label{fig:l15-timing}
\end{figure}

Multiplexing reduces the address pins from $r+c$ to $\max(r,c)$, half as many
when rows and columns are equally many, which makes the package
smaller and cheaper and reduces the wiring between the chip and the memory
controller that drives it (\cref{sec:l15-connection}).  A chip with $2^{28}$
cells in \num{16384} rows and \num{16384} columns, for example, needs 14
address pins instead of~28.

\section{Fast page mode}
\label{sec:l15-fpm}

In the sequence of \cref{fig:l15-timing}a every access opens a row, uses one
column and closes the row again, even when the next access is to the same
row.\source{Lesson 15, videos 11--12.}  Once a row has been opened, however,
all its bits are in the latch, and reaching another bit of the row requires
only a new column address.

\begin{definition}[Fast page mode]
  In \term{fast page mode}[ファストページモード] a row, once opened with RAS,
  stays open while several column addresses are applied with successive CAS
  pulses, each reading or writing one bit of the row.  The row is written back
  only once, when RAS is finally deasserted (\cref{fig:l15-timing}b).
\end{definition}

\caution{This \enquote{page} is a row of the chip, not a virtual-memory page
(Lesson 13).}

Fast page mode pays for the expensive steps---selecting and sensing the row,
and writing it back---once for a whole group of accesses, and for each access
only the quick selection of a column.\sidenote{Fast page mode was the first of
a line of improvements: extended data out (EDO) kept the data valid while the
next column address was applied, synchronous DRAM (SDRAM, first shipped in
1992) replaced the strobes with a clock and a command protocol, and DDR added
a transfer on both clock edges.  Inside the chip the row-and-column structure
is unchanged.}  Accesses to nearby addresses, which
spatial locality (Lesson 12) makes common, usually differ only in the lower
bits of the address and therefore share a row.  The controller that drives RAS and CAS decides
how long a row stays open.  It can keep the row open as long as the following
accesses use the same row, and it must close the row before it accesses
another row or refreshes one.  The savings therefore depend on how many
consecutive accesses share a row.

\begin{example}
  \label{ex:l15-fpm}
  In a DRAM chip, opening a row (RAS, decoding and sensing) takes
  \SI{30}{\nano\second}, selecting a column with CAS \SI{5}{\nano\second}, and
  closing the row (write-back) \SI{10}{\nano\second}.  Five accesses arrive in
  the order of \cref{tab:l15-fpm}, where the pairs give (row, column).
  Without fast page mode every access opens and closes its row, which takes
  $5\times\SI{45}{\nano\second} = \SI{225}{\nano\second}$.  With fast page
  mode, rows are opened and closed three times, and the table adds up to
  \SI{145}{\nano\second}.  Had the access to row~22 come last, the four
  accesses to row~7 would have shared a single opening, for a total of
  \SI{105}{\nano\second}.\caution{Module timings such as 22-22-22 are bus
  clocks, not nanoseconds: on DDR4-3200 opening a row, reading a column and
  closing a row each take \SI{13.75}{\nano\second}.  JEDEC speed bins.}
\end{example}

\begin{table}[htbp]
  \centering\small
  \caption{The accesses of \cref{ex:l15-fpm} in fast page mode.}
  \label{tab:l15-fpm}
  \begin{tabular}{@{}llr@{}}
    \toprule
    Access         & Steps                                      & Time (ns) \\
    \midrule
    read (7, 12)   & open row 7, CAS column 12                  & 35 \\
    read (7, 13)   & CAS column 13                              & 5 \\
    write (7, 14)  & CAS column 14 with write enable            & 5 \\
    read (22, 3)   & close row 7, open row 22, CAS column 3     & 45 \\
    read (7, 15)   & close row 22, open row 7, CAS column 15    & 45 \\
    ---            & close row 7                                & 10 \\
    \midrule
    Total          &                                            & 145 \\
    \bottomrule
  \end{tabular}
\end{table}

\section{Connecting DRAM to the processor}
\label{sec:l15-connection}

The time the processor waits for main memory is not just the access time of
the DRAM chips.\source{Lesson 15, videos 13--14.}  Between the two lies a path
through other chips and buses, shown in its traditional form in
\cref{fig:l15-connection}a.

\begin{definition}[Front-side bus and northbridge]
  The \term{front-side bus}[フロントサイドバス]\index{FSB|see{front-side bus}}
  (FSB) connects the processor to the \term{northbridge}[ノースブリッジ], a
  chip that links the processor with main memory, with fast devices such as
  the graphics adapter, and with the southbridge, to which the slower I/O
  devices are attached (Lesson 16).
\end{definition}

The northbridge contains the \term{memory controller}[メモリコントローラ].  It
translates each memory access that arrives over the FSB into the signals of
the DRAM chips: it splits the address into the row and the column address and
drives RAS, CAS and write enable over the memory bus.  A memory access on a
cache miss therefore travels from the processor over the FSB to the
northbridge and from there over the memory bus to the DRAM chips, and the data
returns the same way.  Each of these steps adds latency, and the FSB carries
all other traffic between the processor and the devices behind the northbridge
as well.\margin{For scale: a 1333~MT/s front-side bus 64 bits wide carries
about 10.7~GB/s for all traffic together, while one DDR4-3200 channel delivers
25.6~GB/s.}

\begin{figure}[htbp]
  \centering
  % Path from the processor to DRAM.  (a) Traditional organization: the
% processor reaches the memory controller in the northbridge over the
% front-side bus; the northbridge drives the memory bus.  (b) Integrated
% memory controller: the memory bus is attached to the processor itself.
\begin{tikzpicture}[x=1cm,y=1cm, font=\footnotesize\sffamily, >={Stealth[length=1.8mm]},
  box/.style={draw, fill=ln-box, align=center, inner sep=3pt},
  mc/.style={draw, fill=ln-accent!15, align=center, inner sep=2pt, font=\scriptsize\sffamily},
  lab/.style={font=\scriptsize\sffamily, text=ln-muted}]
  % \lndimms{x}{y}: three DRAM modules centred at (x,y)
  \def\lndimms#1#2{%
    \foreach \k in {0,1,2} {
      \draw[fill=white] (#1-0.75+0.12*\k,#2-0.45+0.14*\k) rectangle (#1+0.55+0.12*\k,#2+0.25+0.14*\k);
    }
    \node[font=\scriptsize\sffamily] at (#1+0.04,#2+0.08) {DRAM};
  }
  % ================================================= (a)
  \node[anchor=west] at (0,1.9) {\iflangja{(a) ノースブリッジ経由}{(a) through the northbridge}};
  \node[box, minimum width=1.9cm, minimum height=1.0cm] (cpu) at (1.0,0) {\iflangja{プロセッサ}{processor}};
  \node[box, minimum width=2.3cm, minimum height=1.3cm] (nb) at (5.1,0) {};
  \node[anchor=north] at (5.1,0.6) {\iflangja{ノースブリッジ}{northbridge}};
  \node[mc] at (5.1,-0.25) {\iflangja{メモリコントローラ}{memory controller}};
  \lndimms{9.6}{0}
  \node[box, font=\scriptsize\sffamily] (gfx) at (5.1,1.5) {\iflangja{グラフィックス}{graphics}};
  \node[box, font=\scriptsize\sffamily] (sb) at (5.1,-1.55) {\iflangja{サウスブリッジ}{southbridge}};
  \draw[<->, very thick, ln-accent] (cpu) -- node[lab, above] {FSB} (nb);
  \draw[<->, very thick, ln-accent] (nb.east) -- node[lab, above] {\iflangja{メモリバス}{memory bus}} (8.8,0);
  \draw[<->] (nb) -- (gfx);
  \draw[<->] (nb) -- (sb);
  \draw[->] (sb.east) -- (7.0,-1.55) node[lab, anchor=west] {I/O};
  % ================================================= (b)
  \node[anchor=west] at (0,-2.55) {\iflangja{(b) メモリコントローラ統合}{(b) integrated memory controller}};
  \node[box, minimum width=4.3cm, minimum height=1.2cm, anchor=west] (cpu2) at (0.05,-3.5) {};
  \node[anchor=west] at (0.2,-3.5) {\iflangja{プロセッサ}{processor}};
  \node[mc, anchor=east] at (4.25,-3.5) {\iflangja{メモリコントローラ}{memory controller}};
  \lndimms{9.6}{-3.5}
  \draw[<->, very thick, ln-accent] (cpu2.east) -- node[lab, above] {\iflangja{メモリバス}{memory bus}} (8.8,-3.5);
\end{tikzpicture}

  \caption{The path between processor and DRAM.  (a)~The memory controller in
    the northbridge is reached over the front-side bus.  (b)~With the memory
    controller integrated into the processor, the memory bus is attached to
    the processor directly.}
  \label{fig:l15-connection}
\end{figure}

Newer processors integrate the memory controller on the processor chip and
attach the memory bus to the processor directly
(\cref{fig:l15-connection}b).\sidenote{AMD's Athlon~64 (2003) was the first
mainstream x86 processor with the memory controller on the processor die, and
Intel followed with Nehalem in 2008.  The rest of the northbridge---the
graphics and PCI~Express links---moved onto the processor as well, leaving a
single chipset for the slow devices.  Each processor chip now has its own
memory, and an access to the memory of another chip crosses the interconnect:
the NUMA machines of Lesson 18.}
The FSB and the northbridge then disappear from the path to memory, which
lowers the latency of every access to main memory.  Even so, a DRAM access
remains far slower than the processor, for the reasons collected in this
lesson, and the caches of Lessons 12--14 remain indispensable.  Lesson 16
turns to storage devices, which are slower still but keep their data without
power.

\begin{keypoint}[Lesson summary]
  RAM is built as SRAM, whose six-transistor cell holds a bit in an inverter
  loop, or as DRAM, whose one-transistor cell (for example a trench cell)
  holds it as charge on a capacitor.  SRAM is fast and expensive and is used
  for caches.  DRAM is dense and cheap and is used for main memory, but its
  reads are destructive and are followed by a write-back, and leakage forces
  every row to be refreshed, so that one row must be refreshed in every
  interval $T/R$.  A DRAM chip selects a row with the row decoder, reads it
  with the sense amplifiers into a latch and selects a bit with the column
  multiplexer.  The row and column addresses share the address pins and are
  marked by RAS and CAS, and fast page mode serves several accesses to one row
  with a single opening.  A memory access traditionally travels over the
  front-side bus to the memory controller in the northbridge and on over the
  memory bus; newer processors integrate the memory controller and drive the
  memory bus directly.
\end{keypoint}

\end{document}
